\subsection{Ultra Terminum}

\bquote{
  Decoupling the underlying consensus from the state-transition has been informally proposed in private for at least two years---Max Kaye was a proponent of such a strategy during the very early days of Ethereum.
}[Dr. Gavin Wood; \href{\citePolkadotLink}{Polkadot Whitepaper} (2016), s2.2]

In essence, \emph{Ultra Terminum} (UT) is not a method to scale a single blockchain, rather a method to scale a \emph{network} of blockchains.

\UT{} differs from existing consensus methods\footnote{
  Whether \emph{Ultra Terminum} is a consensus \emph{method} or not is somewhat unclear.
    On the one hand: a new combination of methods is still a method.
    On the other hand: UT provides a way to modify other consensus methods via the fork rule, and UT doesn't provide a way to run a \emph{single} blockchain.
    This is why I introduced it as a cross-chain consensus \emph{strategy}.
} in that it is a \emph{novel modification} to particular components of preexisting consensus algorithms.
These modifications allow those consensus algorithms to cooperate.
This improves the maximal security and decentralization of the network, whilst also providing a foundation for scalability.

At the core of \emph{Ultra Terminum} is a new method for sharing security: \emph{Proof of Reflection} (\autoref{sec:proof-of-reflection}).
This technique (which works \emph{in conjunction} with PoW, PoS, etc) allows the incremental construction of complex blockchain networks with powerful scaling properties.
Those blockchains use the technique \emph{Proof of Reflection} to \emph{mutually secure each other}, forming something called \emph{the simplex} (\autoref{sec:the-simplex}).
Blockchains at this level are called \emph{simplex-chains} and are able to host \emph{dapp-chains} --- chains that are dedicated to a dapp, and they can be application specific (e.g., a DEX) or general (e.g., a chain hosting an instance of the EVM).
The vital properties provided by \UT{} allow heterogenous chains to form a cohesive network and cooperate in securing each other.
\emph{Proof of Reflection} is how \emph{Ultra Terminum} (excluding nested chains) scales with order $O(c^2)$, and is the basis for unbounded $O(n)$ scaling.
UT's higher-order scaling configurations, $O(c^3)$ and $O(c^4)$, require dapp-chains: chains that are application-specific and which inherit security properties from the foundational structure.
As UT is primarily a method of \emph{structuring} a blockchain network, the scalability configurations mentioned herein \emph{do not include \textbf{layer 2} methods}.
That means that \emph{layer 2} techniques (e.g., state/payment channels, ZK/optimistic rollups) can be implemented \emph{on top} of UT.

UT's novel structure means that confirmation times within UT are \ul{\emph{sub-constant}} and of order $O(c^{-1})$ (i.e., the confirmation rate is $O(c)$).
This means that, as the network grows and $c$ increases, confirmation times in UT will approach 0.
This is an improvement over existing architectures, which are, ideally and at best, $O(1)$.

Practically, there is little difference between a single blockchain and the simplex to the user.
\emph{Proof of Reflection} allows each simplex-chain, and each hosted dapp-chain, to do efficient SPV proofs regarding state-entries on other simplex-chains or dapp-chains.
This allows for the creation of rich cross-chain functionality, such as the ability to use a single network-level token across all simplex-chains and dapp-chains.

Although SPV proofs can be hundreds of bytes long, UT is structured so that their length does not significantly impact scalability.
%\footnote{
%    Also, verkle trees can replace merkle trees --- this reduces proof costs (approximately) by $5 \text{ to } 10 \times$.
%}
UT does this by prioritizing a clean and simple protocol for simplex-chains, and allowing dapp-chains the freedom to implement whatever state- and transaction-protocols are suitable for the application in question.
This means that complex or inefficient data-structures for a dapp-chain's state do not impact the network globally, preventing bottlenecks in foundational components.

Finally, a UT simplex can be modified to enable a \emph{tiling of simplexes} (\autoref{sec:tiling}).
This technique \emph{removes the upper bound} on a UT network's growth, without sacrificing security of the network.
